%==============================================================================
\section{Act II: The Brain (\texttt{System.h/.cc})}
%==============================================================================

The \texttt{System} class is the \textbf{orchestrator}. It doesn't do any computation itself, but it creates and coordinates all the components.

\subsection{What Does the System Own?}

Looking at \texttt{System.h}, we see the System holds pointers to everything:

\begin{lstlisting}[caption={System's private members (simplified)}]
private:
    eSensor mSensor;              // MONOCULAR, STEREO, RGBD, etc.
    
    // The knowledge base
    ORBVocabulary* mpVocabulary;  // Visual vocabulary for recognition
    KeyFrameDatabase* mpKeyFrameDatabase;  // For loop detection
    Atlas* mpAtlas;               // Collection of all maps
    
    // The workers
    Tracking* mpTracker;          // Estimates camera pose
    LocalMapping* mpLocalMapper;  // Builds the map
    LoopClosing* mpLoopCloser;    // Fixes drift
    
    // Visualization
    Viewer* mpViewer;
    FrameDrawer* mpFrameDrawer;
    MapDrawer* mpMapDrawer;
    
    // Threads
    std::thread* mptLocalMapping;
    std::thread* mptLoopClosing;
    std::thread* mptViewer;
\end{lstlisting}

\subsection{The System Constructor: Birth of SLAM}

When you call \texttt{System(vocabPath, settingsPath, MONOCULAR, true)}, this happens:

\begin{enumerate}
    \item \textbf{Load Vocabulary}: Read the ORB dictionary from disk (~100MB, takes a few seconds)
    \item \textbf{Create Atlas}: The ``library'' that will store all maps and map points
    \item \textbf{Create KeyFrameDatabase}: For finding similar places (loop closure)
    \item \textbf{Create Drawers}: For visualization
    \item \textbf{Create Tracking}: The main worker (we'll explore this in the next section)
    \item \textbf{Create LocalMapping}: Spawn a background thread
    \item \textbf{Create LoopClosing}: Spawn another background thread
    \item \textbf{Create Viewer}: Spawn the visualization thread
\end{enumerate}

\subsection{The Sensor Types}

The System supports multiple sensor configurations:

\begin{lstlisting}[caption={Sensor types enum}]
enum eSensor {
    MONOCULAR = 0,     // Single camera
    STEREO = 1,        // Two cameras (left/right)
    RGBD = 2,          // RGB + Depth camera (Kinect, RealSense)
    IMU_MONOCULAR = 3, // Mono + IMU
    IMU_STEREO = 4,    // Stereo + IMU
    IMU_RGBD = 5       // RGBD + IMU
};
\end{lstlisting}

\subsection{The Public Interface}

The System exposes a simple interface to the outside world:

\begin{lstlisting}[caption={System's main interface}]
// Feed images to the system (returns camera pose)
Sophus::SE3f TrackMonocular(const cv::Mat &im, const double &timestamp);
Sophus::SE3f TrackStereo(const cv::Mat &imLeft, const cv::Mat &imRight, ...);
Sophus::SE3f TrackRGBD(const cv::Mat &im, const cv::Mat &depthmap, ...);

// Control
void ActivateLocalizationMode();  // Stop mapping, only track
void DeactivateLocalizationMode();
void Reset();
void Shutdown();

// Output
void SaveTrajectoryTUM(const string &filename);
void SaveTrajectoryEuRoC(const string &filename);
\end{lstlisting}

\textbf{Key insight:} The System class follows the \textit{Facade} design pattern. It hides all the complexity behind a simple interface. You don't need to know about Frames, KeyFrames, or MapPoints to use it.

\subsection{What Happens in \texttt{TrackMonocular}?}

When you call \texttt{SLAM.TrackMonocular(im, timestamp)}:

\begin{lstlisting}[caption={TrackMonocular simplified flow}]
Sophus::SE3f System::TrackMonocular(const cv::Mat &im, const double &timestamp)
{
    // 1. Convert to grayscale if needed
    cv::Mat imGray;
    if(im.channels() == 3)
        cv::cvtColor(im, imGray, cv::COLOR_BGR2GRAY);
    else
        imGray = im;
    
    // 2. Delegate to Tracking (the real worker!)
    Sophus::SE3f Tcw = mpTracker->GrabImageMonocular(imGray, timestamp, "");
    
    // 3. Update state for external queries
    mTrackingState = mpTracker->mState;
    mTrackedMapPoints = mpTracker->mCurrentFrame.mvpMapPoints;
    
    return Tcw;  // Return the camera pose
}
\end{lstlisting}

So the System is essentially a thin wrapper. The real work happens in \texttt{Tracking}.

\subsection{The Multi-Threading Model}

ORB-SLAM3 uses a producer-consumer pattern:

\begin{itemize}
    \item \textbf{Tracking (Main Thread)}: Produces KeyFrames, consumes images
    \item \textbf{LocalMapping (Thread 1)}: Consumes KeyFrames, produces MapPoints
    \item \textbf{LoopClosing (Thread 2)}: Consumes KeyFrames, produces corrections
    \item \textbf{Viewer (Thread 3)}: Consumes map state, produces visualization
\end{itemize}

Communication happens through:
\begin{itemize}
    \item Shared pointers to KeyFrames and MapPoints
    \item Thread-safe queues
    \item Mutex-protected data access
\end{itemize}
